Let's see how we can generalize what we have discussed so far. First of all, we are gonna to describe some useful definitions, as follows.
\begin{definition}[title = Instance Space $X$]
    We assume to work within an \textbf{instance space} $\boldsymbol{X}$:
    \begin{itemize}
        \renewcommand{\labelitemi}{-}
        \item $X$ is the set of instances of objects the learner wants to classify.
        \item Data from the instance spaces are generated through a distribution $\boldsymbol{D}$.
        \item The instance space $X$ is the set where training data and testing data belong.
    \end{itemize} \vspace{7pt}

    Instance space can be seen as the set of \textbf{raw data}, that follows an unknown distribution $\boldsymbol{D}$ and does not include labels at all. \vspace{7pt}

    Note: labelled data will be provided to our learning algorithm by an \textbf{oracle}, formally described as $EX(c, \boldsymbol{D})$.
\end{definition} 
\begin{definition}[title = Concept]
    \textbf{Concepts} are subsets of $X$, thought as properties of data or objects. A concept is an arbitraly subset of the instance space $X$. \vspace{7pt}

    In the previous example about rectangles, concepts are arbitrary subsets of the instance space $X = R^{2}_{[0,1]}$, since we take pairs of coordinates. For 
    instance:
    \begin{itemize}
        \renewcommand{\labelitemi}{-}
        \item Rectangle has parallel axes.
        \item Red points are inside the rectangle.
        \item Blue points are outside the rectangle.
    \end{itemize}
    However, these concepts are not provided as input to the learning algorithm: it should derive these concepts while learning patterns from training data.
\end{definition}
\begin{definition}[title = Concept class]
    A \textbf{concept class} $C$ is a collection of concepts, namely a subset of $P(X)$. These are the concepts we would like to learn during the training 
    phase. \vspace{7pt}

    For instance, the concept class in the previous example that we learnt is the one of rectangles whose sides are parallel to the axes. \vspace{7pt}

    Note: the learning algorithm did not know that the classifier would be a rectangle whose sides are parallel to axes, only during the training phase 
    it recognized this behavior from data. \vspace{7pt}

    The \textbf{target concept}, namely $c \in C$, is the concept, so the property, that the learner wants to build a classifier for (\textit{boolean function to 
    recognize when a point is inside or outside the rectangle}).
\end{definition}